% SI numbering: first-citation order (tools/si_numbering.py, round 6, 2026-10-05)
% results_a.tex : 결과 R1-R4(국문판). 영문판 paper/manuscript/en/sections/results_a.tex(2026-10-05 2차 수정판)를
% 문장 단위로 옮겼다. 수치, 인용 키, 그림·보충 표 번호, 자리표시 문자열은 영문판과 같다.
% 신뢰구간은 [하한, 상한]으로 적는다(값은 영문판의 'a to b' 와 같다).

\section{결과}\label{sec:results}

\subsection{평가 설계와 검증 사다리}\label{sec:results-design}

앞에서 설명한 홀드아웃 절차(그림~1c,d)에서 모든 방법을 같은 대상 라벨을 쓰는 두 Stefan 기준선과 비교하였다(그림~2a,b). 하나는 원천 셀에서 계수 $E_0$를 추정한 원천 계수 Stefan 식(표~1)이고, 다른 하나는 대상 라벨로 계수를 다시 맞추되 $E_0$ 쪽으로 수축한 재보정 Stefan 식($\kappa = 10$)이다\cite{Steyerberg2004}.
4지역 평균은 레나델타, 캐나다, 러시아 서부, 러시아 동부에 같은 가중을 주며, 학습에서 뺀 참조 대상인 알래스카는 넣지 않는다.
구간은 블록 등가중이라고 적지 않으면 셀 가중이다.

다섯 지역을 같은 가중으로 평균하면 직접 ML의 RMSE는 무작위 셀 분할의 22.9~cm에서 지역 홀드아웃의 33.3~cm로 커졌다(차이 10.3~cm, 95\% CI [7.3, 12.8], 블록 등가중 6.1~cm [4.7, 7.6]). 같은 조건에서 학습 셀에 최소제곱으로 계수를 맞춘 Stefan 식은 26.8--27.0~cm에 머물렀다(그림~1e).
10.3~cm 증가는 확인적 결과가 아니라 효과 크기로 보고하며, 등록 대비인 500~km 완충 군집 홀드아웃과 무작위 셀의 차이는 8.2~cm였다(95\% CI [5.6, 10.4]; 보충 표~S1).
무작위 셀 분할에서 $k$겹 최근접 거리 정합(k-fold nearest-neighbour distance matching, kNNDM) 분할\cite{Linnenbrink2024}로 가면 직접 ML의 RMSE는 알래스카, 레나델타, 캐나다에서 5.76~cm, 8.17~cm, 14.20~cm 늘었고, 학습 셀에 맞춘 Stefan 식의 RMSE는 0.29~cm, 0.51~cm, 4.28~cm 늘었다(셀 가중; 서술, 앞선 결과를 본 뒤 설계; 그림~1e; 보충 표~S2).
확률 표본에서는 무작위 검증이 지도 정확도의 타당한 추정이라는 견해가 있으나\cite{Wadoux2021}, 이 연구의 라벨은 군집된 비확률 표본이므로 아래의 오차는 학습에서 뺀 블록에서의 예측 오차다.

\subsection{대상 라벨 없는 전이}\label{sec:results-label0}

대상 라벨이 없을 때 Stefan 식의 유사라벨은 섞은 유사라벨보다 RMSE를 1.6~cm 낮추었다(95\% CI [1.0, 2.0], 블록 등가중 [0.6, 1.5], 4지역, Holm 보정 \textit{P} = 0.001). 사후 집계에서 원천 계수 Stefan 식보다 오차가 2~cm 넘게 커진 대상은 30개 가운데 물리 유사라벨 증강에서 2개, 직접 ML에서 18개였다(그림~3b; 그림~4a,b).
두 상수 위약에 대한 감소는 1.7~cm와 3.6~cm였고, $\sqrt{\mathrm{TDD}}$에 선형인 위약에 대해서는 0.48~cm였다(등록 판정: 혼재).
2~cm 문턱은 전이 결과를 본 뒤 정하였고, 30개 대상에는 하위 지역 또는 라벨 확충판 23개가 들어 있으며, 구간에 근거한 집계도 같은 순서를 보였다(사후; 보충 표~S3).

주 구간과 공통 재표집 구간의 판정이 일치한 학습기 5종에서 직접 ML은 원천 계수 Stefan 식보다 4지역 평균 오차가 높았다(그림~4c).
주 학습기인 CatBoost에서 차이는 2.2~cm(95\% CI [1.1, 3.4], Holm 보정 \textit{P} = 0.023)였고, CatBoost의 다른 설정 2종과 TabPFN v2, TabICL v2에서는 0.85--1.9~cm였다.
이 다섯 판정은 분할 독립 가정에 의존하며, TabPFN v2의 판정은 보정 전에만 유의하였다(Holm 보정 \textit{P} = 0.117).
원천 계수 앵커에 낮은 가중($\lambda = 0.25$)을 준 잔차 모델은 학습기 3종에서 그 Stefan 식과 ±0.5~cm 안에서 동등하였다(분할 독립 가정에 의존).

알래스카를 학습에서 뺐을 때 물리 유사라벨은 원천 계수 모델보다 오차가 1.8~cm 높았다(95\% CI [0.4, 2.6]).
원천 계수가 지역 값과 멀리 떨어진 티베트 고원(그림~1b)에서는 직접 ML의 오차가 60.0~cm 낮았다(비맹검이고 교차 환경인 비교에서 예상된 결과).
라벨 10개에서 직접 TabICL v2는 4지역 평균 오차가 2.9~cm 낮았으나(95\% CI [1.1, 3.6]; 분할 독립 가정 의존), 알래스카를 포함한 3지역 평균에서는 6.0~cm 높았다.

같은 블록에서 제품 학습 지점으로부터 5~km 또는 25~km 안의 블록을 가린 뒤(등록 이탈; 앞선 결과를 본 뒤 설계), 라벨 없는 기존 CCI v5 지도는 원천 계수 Stefan 식보다 오차가 23.79~cm 높았고(Holm 보정 \textit{P} = 0.0006; 레나델타, 캐나다, 러시아 동부에서 높음), 잔차 ML은 같은 라벨 10개와 전량으로 재보정한 CCI v5보다 오차가 8.37~cm와 7.51~cm 낮았다(CCI v5 대비는 비맹검 부분 포함). 대상 지역 안의 CALM 지점도 일부 학습에 쓴 Wei v2 지도(제품 결측 대체 47.4\%)에서는 차이를 확인하지 못했다(ALT 정의가 서로 다름; 보충 표~S2).

\subsection{라벨 3--10개의 계수 재보정}\label{sec:results-recal}

대상 라벨 10개로 Stefan 계수를 재보정하면 원천 계수보다 4지역 평균 RMSE가 2.5~cm 낮아졌다(95\% CI [1.9, 3.0], 블록 등가중 [1.9, 3.3], Holm 보정 \textit{P} = 0.001). 오차가 낮아진 지역은 4곳 가운데 1곳(러시아 서부, 12.0~cm)이었고, 캐나다에서는 오차가 높아졌으며(2.3~cm, 95\% CI [0.7, 2.9]), 레나델타와 러시아 동부에서는 차이를 확인하지 못했다(지역 중앙값 변화 −0.06~cm, 러시아 서부 제외 +0.72~cm; 그림~2c--f; 그림~3a,d).
재보정 앵커에 잔차 ML을 더하면 RMSE가 추가로 −0.18~cm 달라졌다(95\% CI [−0.53, −0.01], Holm 보정 \textit{P} = 0.136). 이 차이는 보정 전에만 유의하고, 크기가 0.5~cm 미만이며, 분할 독립 가정에 의존한다.

라벨 전량에서 라벨을 쓰는 보정(다음 절의 선정 규칙)의 오차 감소는 라벨 10개로 잰 원천 계수 편향과 순위상관을 보였다(Spearman $\rho = 0.38$, 95\% CI [0.17, 0.55], 28대상, 단측 \textit{P} = 0.0002; 라벨 격자 결과를 본 뒤 설계).
편향의 부호는 감소와 관계가 없었다($\rho = -0.01$; 그림~5a,b; 보충 표~S4).
사후 분석에서 이 편향의 크기와의 관계는 재보정 이득, 선정 규칙이 더한 몫, 재보정을 넘어선 ML 몫의 어느 것에서도 확인하지 못했다(마지막 몫은 대상 고정 보조 구간에서만 양이었고, 약한 쪽 범주로 적었다)(그림~5d; 보충 표~S2).

라벨 10개 이하에서 잔차 구조는 Stefan 출력을 입력으로 받는 ML(물리 입력 ML)보다 4지역 평균 오차가 낮았다. 차이는 라벨 0개에서 2.7~cm, 라벨 10개에서 3.7~cm(95\% CI [2.8, 5.2], Holm 보정 \textit{P} = 0.001; 라벨 10개에서 4지역 가운데 2지역의 오차가 낮았고, 지역 중앙값 변화 −2.2~cm, 러시아 서부 제외 −0.32~cm)였다(그림~3c).
재보정 Stefan 값을 물리 입력으로 주어도 잔차 구조의 오차는 2.5~cm 낮았으나 캐나다에서는 3.1~cm 높았고, 이 지역 판정은 분할 독립 가정에 의존한다(보충 표~S5).

잔차 ML에 물리 유사라벨을 더한 결과는 라벨 10개에서 ±0.5~cm 안의 동등이었고, 라벨 40개의 레나델타와 캐나다에서는 오차가 0.34~cm 높았다(0.5~cm 미만; 보충 표~S5).
물리식과 제품을 라벨 가중으로 적층한 결과는 라벨 10개에서 재보정이나 잔차 ML과의 차이를 확인하지 못했고, 이 라벨 수에서 적층은 추출의 36\%에서 재보정으로 대체되었다(앞선 결과를 본 뒤 설계, 비맹검 부분 포함; 보충 표~S2).

\subsection{라벨 수십--수백 개의 방법 선정}\label{sec:results-selection}

대상 라벨 안에서 5겹 블록 교차검증으로 방법을 고르면 고정 잔차 레시피($\lambda = 0.25$)보다 라벨 160개에서 RMSE가 1.3~cm 낮아졌고(95\% CI [0.7, 1.7], 레나델타와 캐나다, 공통 재표집에서도 같음), 라벨 40개와 160개에서 비열등이었다(한계 0.5~cm). 두 결과 모두 캐나다에서 나왔으며, 이 규칙은 라벨 격자 결과를 본 뒤 설계하였다(그림~5c; 그림~6d).
라벨 40개에서의 감소 0.87~cm는 분할 독립 가정에 의존한다.
레나델타에서는 차이를 확인하지 못했다.
주 4지역의 라벨 10개에서는 비열등을 확인하지 못했고, 레나델타에서 규칙의 오차가 0.70~cm 높았다. 라벨 전량에서의 비열등은 분할 독립 가정에 의존한다.
규칙은 재보정 Stefan 식보다 오차가 낮은 대상의 수를 늘리지 않았다(보충 표~S4).

캐나다에서는 원천 계수 Stefan 식과의 차이를 확인하지 못했다(보충 표~S4).
알래스카를 학습에서 뺐을 때 규칙은 라벨 10, 40, 160개에서 원천 계수 모델보다 오차가 높았다(서술적 곡선 대비; 보충 표~S4).

라벨 40개와 160개에서 레나델타와 캐나다의 잔차 ML은 물리 입력 ML보다 오차가 1.9--2.4~cm 높았다(2.4~cm 대비는 분할 독립 가정에 의존)(그림~3c).
이 차이는 캐나다에서 생겼고, 알래스카를 학습에서 뺐을 때는 잔차 ML의 오차가 낮았다.
이 라벨 수에서 잔차 ML과 더 강한 재보정 기준선(토양 보정 Stefan 식을 $\kappa = 10$으로 재보정한 것) 사이의 차이는 확인하지 못했다.

원천 지역 하나 제외 교차검증으로 신경망을 조정한 결과는 학습기마다 달랐고, 학습기 전반에 대한 문장을 쓸 등록 조건을 채운 학습기는 없었으며, 이 조정은 대상 라벨 안의 선택과 비교하지 않았다(보충 표~S4).
같은 지역을 다시 무작위로 나눈 분할에서(재사용 지역의 재검정, 앞선 결과를 본 뒤 설계, 비맹검 부분 포함) 순차 워크플로는 무작위 배치와 고정 잔차 레시피보다 라벨 160개에서 오차가 0.93~cm 낮았고(레나델타와 캐나다; 캐나다 −2.41~cm, 레나델타 +0.54~cm 미결정), 라벨 40개에서는 차이를 확인하지 못했다(최소 검출 효과 약 0.5~cm; 레나델타 +0.21~cm). 워크플로의 배치 단계가 무작위 추출로 바뀌었으므로 이 대비는 같은 라벨에서 방법 규칙과 고정 레시피를 비교한 것이다.
같은 분할에서 워크플로는 라벨 10개, 40개, 160개에서 재보정 Stefan 식에 비열등이었고(라벨 10개는 5지역 평균, 40개와 160개는 레나델타와 캐나다 평균) 라벨 40개와 160개에서는 오차가 1.35~cm와 1.51~cm 낮았으나, 이는 풀 평균에 대한 것이다. 러시아 동부의 라벨 10개(+1.34~cm), 선택 계열인 알래스카의 라벨 40개와 160개(+1.51~cm와 +2.34~cm), 레나델타의 라벨 160개(+0.53~cm)에서는 0.5~cm 넘게 나빴다(그림~6e; 보충 표~S2).
지역 안의 보조 팔에서 규칙은 라벨 200개에서 잔차 ML과 동등하였고, 라벨 500개와 1000개에서는 차이를 확인하지 못했다.
적층은 라벨 40개와 160개에서도 차이를 확인하지 못했다(레나델타와 캐나다; 앞선 결과를 본 뒤 설계, 비맹검 부분 포함; 보충 표~S2).
